% Supplementary Material — standalone \input{} file
% Requires: amsmath, amssymb, booktabs, array, float, enumitem packages in parent

\section{Extended Derivations}
\label{app:derivations}

\subsection{Geodesic distance on the Grassmannian}
\label{app:geodesic}

The Grassmannian $\Gr(k,d)$ admits a canonical Riemannian metric inherited
from the ambient space of $d \times k$ orthonormal frames. Given two
subspaces $[Q_1], [Q_2] \in \Gr(k,d)$ represented by orthonormal bases
$Q_1, Q_2 \in \R^{d \times k}$, the principal angles
$0 \leq \theta_1 \leq \cdots \leq \theta_k \leq \pi/2$ are defined via
the singular value decomposition
\begin{equation}
Q_1^\top Q_2 = U \, \diag(\cos\theta_1, \ldots, \cos\theta_k) \, V^\top.
\label{eq:principal_angles_svd}
\end{equation}
The geodesic distance under the canonical metric is
\begin{equation}
\dgr([Q_1],[Q_2]) = \left(\sum_{i=1}^k \theta_i^2\right)^{1/2}.
\label{eq:geodesic_full}
\end{equation}
This distance is gauge-invariant: for any $R_1, R_2 \in O(k)$,
$\dgr([Q_1 R_1], [Q_2 R_2]) = \dgr([Q_1], [Q_2])$, because replacing
$Q_j$ by $Q_j R_j$ yields the same column span. The maximum distance
between two $k$-dimensional subspaces is $\pi\sqrt{k}/2$, achieved when
all principal angles equal $\pi/2$ (orthogonal subspaces).

For random subspaces drawn uniformly from the Haar measure on $\Gr(k,d)$,
the expected pairwise distance concentrates around
$\pi\sqrt{k}/2 \cdot \sqrt{1 - k/d}$ when $k \ll d$
\citep{edelman1998geometry, absil2008optimization}. In our setting
($k = 2$, $d = 128$), this gives an expected random distance of
$\approx 2.20$ radians, consistent with the observed inter-seed distances
of $[1.93, 2.19]$ (\S\ref{sec:gauge_freedom}).

\subsection{Equivariance test statistic}
\label{app:equiv_test}

For an operation with a cyclic group action $a \mapsto a + g \pmod{p}$,
we test whether the DAS subspace respects the group symmetry. Let
$z(a,b) = Q^\top h(a,b) \in \R^k$ be the DAS projection. The equivariance
condition requires that the group action on the input induces a
consistent rotation in the projected space:
\begin{equation}
z(a + g, b) = R_g \, z(a, b) + t_g \quad \forall\, a, b, g
\label{eq:equiv_condition}
\end{equation}
for some rotation $R_g \in O(k)$ and translation $t_g$ depending only on
$g$, not on $(a,b)$.

For $k = 2$ and cyclic groups $\Z_p$, we parameterize $R_g$ as a rotation by
angle $\phi_g$. We estimate $\phi_g$ for each $g$ by computing the
median angular displacement across all $(a,b)$ pairs:
\begin{equation}
\hat\phi_g = \mathrm{median}_{(a,b)} \;\angle\bigl(z(a{+}g, b) - \bar{z},\;
z(a, b) - \bar{z}\bigr)
\label{eq:angle_estimate}
\end{equation}
where $\bar{z}$ is the grand mean of all projections and $\angle(\cdot, \cdot)$
denotes the signed angle. An input $(a,b)$ is counted as equivariant for shift
$g$ if $|\angle - \hat\phi_g| < \varepsilon$ (we use $\varepsilon = 2\pi/p$,
one ``slot'' of angular resolution). The reported equivariance fraction is
\begin{equation}
\text{Equiv} = \frac{1}{|\mathcal{T}|} \sum_{(a,b,g) \in \mathcal{T}}
\mathbf{1}\bigl[|\angle(z(a{+}g,b), z(a,b)) - \hat\phi_g| < \varepsilon\bigr]
\label{eq:equiv_fraction}
\end{equation}
where $\mathcal{T}$ is the set of all valid (input, shift) triples on the
test set. Random orthogonal subspaces of the same dimension achieve
equivariance $< 1\%$ across all operations tested, confirming that the
metric has strong discriminative power.

\subsection{Diversity ratio}
\label{app:diversity_ratio}

The intervention diversity ratio $\rho$ (Eq.~\ref{eq:diversity_ratio})
measures whether an encoder-decoder intervention preserves input-specific
information or collapses to a class-conditional lookup table.

For a set of intervened activations $\{h'_i\}$ sharing source label $y_s$,
the within-class standard deviation is
$\sigma_{y_s} = \bigl(\frac{1}{n_{y_s}} \sum_i \|h'_i - \bar{h}'_{y_s}\|^2
\bigr)^{1/2}$
where $\bar{h}'_{y_s}$ is the mean intervened activation for source label $y_s$.
The denominator uses the same statistic computed on natural (non-intervened)
activations sharing the same label. The ratio $\rho = \mathbb{E}_{y_s}[\sigma_{y_s}^{\text{iv}}]
/ \mathbb{E}_{y_s}[\sigma_{y_s}^{\text{nat}}]$ ranges from 0 (complete collapse: all
interventions for the same source produce the same activation) to
$\approx 1$ (intervention preserves the natural within-class variability).

Values of $\rho \gg 1$ are possible in principle (the intervention adds noise),
but in practice indicate that the decoder is amplifying base-specific variation.
The observed NL-DAS values ($\rho \approx 0.05$ on IOI) confirm near-complete
collapse; the Structured pi-SAE values ($\rho \approx 0.89$) confirm
preservation of nuisance information.


\section{Hyperparameter Tables}
\label{app:hyperparameters}

\subsection{Grokking model training}

\begin{table}[H]
\centering
\small
\setlength{\tabcolsep}{4pt}
\begin{tabular}{ll}
\toprule
\textbf{Parameter} & \textbf{Value} \\
\midrule
Architecture & 1-layer transformer \\
$d_\text{model}$ & 128 \\
$n_\text{heads}$ & 4 \\
$d_\text{head}$ & 32 \\
$d_\text{MLP}$ & 512 \\
Optimizer & AdamW \\
Learning rate & $10^{-3}$ \\
Weight decay & 1.0 \\
Batch size & Full batch \\
Train/test split & 30\% / 70\% \\
Training epochs & 25{,}000--80{,}000 (operation-dependent) \\
Deeper models & 1.5$\times$ epochs for 2L and 4L \\
Grokking criterion & Test cross-entropy $< 0.1$ \\
Moduli tested & $p \in \{53, 97, 113, 211\}$ \\
Depths tested & $n_\text{layers} \in \{1, 2, 4\}$ \\
\bottomrule
\end{tabular}
\caption{Grokking model training hyperparameters. All operations use the same
architecture and optimizer; only epoch count varies.}
\label{tab:hp_grokking}
\end{table}

\subsection{DAS fitting}

\begin{table}[H]
\centering
\small
\begin{tabular}{ll}
\toprule
\textbf{Parameter} & \textbf{Value} \\
\midrule
Optimizer & Adam \\
Learning rate & $10^{-3}$ \\
Gradient steps & 200 \\
Batch size & 16 interchange pairs \\
$k$ values swept & $\{2, 4, 6, 8, 10, 12, 16, 20, 24, 32\}$ \\
Initialization & Random orthonormal $Q \in \R^{d \times k}$ \\
Intervention site & Residual stream (post-MLP, layer 0) \\
Hard-example selection & $y_b \neq y_s$, both correctly classified \\
\bottomrule
\end{tabular}
\caption{DAS fitting hyperparameters. Each $k$ value is an independent
optimization, not a nested subspace.}
\label{tab:hp_das}
\end{table}

\subsection{Structured pi-SAE}

\begin{table}[H]
\centering
\small
\begin{tabular}{ll}
\toprule
\textbf{Parameter} & \textbf{Value} \\
\midrule
Encoder & 2-layer MLP ($d_\text{model} \to 128 \to z$), ReLU \\
Decoder & 2-layer MLP ($z \to 128 \to d_\text{model}$), ReLU \\
$k$ (causal latent dim) & $\{2, 4, 8\}$ \\
$m$ (nuisance latent dim) & 15 \\
Sparsity & $L_1$ on $z_\text{causal}$, expansion $8\times$ \\
Prior & $p(z_\text{causal} \mid y) = \mathcal{N}(\mu_y, \sigma_y^2 I)$ \\
$\beta$ (KL weight) & 1.0 \\
$\alpha$ (classification weight) & 10.0 \\
$\lambda$ ($L_1$ coefficient) & 1.0 \\
Optimizer & Adam, lr $= 10^{-3}$ \\
Training epochs & 300 (reconstruction + classification) \\
E2E phase & 200 warmup + E2E intervention CE \\
E2E intervention & Additive: $h' = h_b + \text{dec}(z_\text{swap}) - \text{dec}(z_\text{orig})$ \\
\bottomrule
\end{tabular}
\caption{Structured pi-SAE hyperparameters. The E2E phase is used for
GPT-2 language tasks; grokking models use reconstruction + classification only.}
\label{tab:hp_pisae}
\end{table}

\subsection{NL-DAS baseline}

\begin{table}[H]
\centering
\small
\begin{tabular}{ll}
\toprule
\textbf{Parameter} & \textbf{Value} \\
\midrule
Encoder $f$ & 2-layer MLP ($d \to 256 \to d$), ReLU \\
Decoder $g$ & 2-layer MLP ($d \to 256 \to d$), ReLU \\
$Q$ & Learned jointly, $\R^{d \times k}$ \\
Training objective & Interchange loss only (no reconstruction) \\
NL-DAS+recon variant & $+ \lambda_\text{recon} \|g(f(h)) - h\|^2$, $\lambda = 0.1$ \\
Optimizer & Adam, lr $= 10^{-3}$ \\
Training steps & 200 \\
\bottomrule
\end{tabular}
\caption{NL-DAS hyperparameters. Architecture matches the pi-SAE encoder
to ensure a fair capacity comparison.}
\label{tab:hp_nldas}
\end{table}


\section{Per-Operation Atlas Data}
\label{app:atlas_data}

Table~\ref{tab:full_atlas} reports the complete atlas measurements for all
14 operations at $p = 113$ (1-layer), including metrics omitted from the
main text for space.

\vspace{1em}
\begin{table}[H]
\centering
\small
\setlength{\tabcolsep}{3pt}
\renewcommand{\arraystretch}{1.15}
\begin{tabular}{@{}lcccccccc@{}}
\toprule
\textbf{Operation} & \textbf{Grok} & \textbf{Loss} & \textbf{IIA}$_{k{=}2}$ &
\textbf{IIA}$_{k{=}4}$ & \textbf{IIA}$_{k{=}8}$ & $k^*$ &
\textbf{Equiv.} & \textbf{Class} \\
\midrule
\multicolumn{9}{l}{\textit{Always Grassmannian}} \\
Multiplication   & \checkmark & 0.000 & 1.000 & 1.000 & 1.000 & 2 & 0.995 & Always \\
Subtraction      & \checkmark & 0.000 & 1.000 & 1.000 & 1.000 & 2 & 1.000 & Always \\
Division         & \checkmark & 0.000 & 1.000 & 1.000 & 1.000 & 2 & 1.000 & Always \\
Bitwise XOR      & \checkmark & 0.003 & 1.000 & 1.000 & 1.000 & 2 & 1.000 & Always \\
Cubic sum        & \checkmark & 0.000 & 1.000 & 1.000 & 1.000 & 2 & 1.000 & Always \\
Sum of squares   & \checkmark & 0.005 & 1.000 & 1.000 & 1.000 & 2 & 0.987 & Always \\
Max              & \checkmark & 0.074 & 0.960 & 1.000 & 1.000 & 2 & 0.957 & Always \\
\midrule
\multicolumn{9}{l}{\textit{Stochastic}} \\
Comp.\ add.\ (orig) & \checkmark & 0.000 & --- & --- & --- & --- & 0.998 & Stochastic \\
Comp.\ add.\ (s42)  & $\times$   & 10.21 & 0.800 & 0.860 & 0.900 & 6 & 0.110 & Stochastic \\
Power (s137)     & \checkmark & 0.088 & 0.940 & 1.000 & 1.000 & 4 & 0.965 & Stochastic \\
Power (s42)      & $\times$   & 1.143 & 0.980 & 1.000 & 1.000 & 2 & 0.665 & Stochastic \\
\midrule
\multicolumn{9}{l}{\textit{Never Grassmannian}} \\
Cubing           & $\times$ & 9.765 & 0.857 & 1.000 & 1.000 & 4 & 0.509 & Never \\
Squaring         & $\times$ & 7.808 & 0.857 & 1.000 & 1.000 & 4 & 0.314 & Never \\
Abs.\ difference & $\times$ & 1.835 & 0.800 & 0.840 & 0.860 & $>$32 & 0.189 & Never \\
Polynomial       & $\times$ & 21.55 & 0.720 & 0.780 & 0.820 & $>$32 & 0.015 & Never \\
Affine           & $\times$ & 26.21 & 0.780 & 0.800 & 0.840 & $>$32 & 0.026 & Never \\
\bottomrule
\end{tabular}
\renewcommand{\arraystretch}{1.0}
\caption{Complete atlas at $p = 113$ (1-layer). IIA is reported at $k \in \{2, 4, 8\}$;
$k^*$ is the smallest $k$ achieving IIA $\geq 0.95$. Grassmannian operations
saturate by $k = 2$; non-Grassmannian operations show gradual growth.
Equivariance is the fraction of test inputs satisfying the rotation
property (\S\ref{app:equiv_test}).}
\label{tab:full_atlas}
\end{table}

Table~\ref{tab:seed_data} reports per-seed DAS results for the 10 seeds
of modular multiplication used in the subspace degeneracy analysis
(\S\ref{sec:gauge_freedom}).

\vspace{1em}
\begin{table}[H]
\centering
\small
\begin{tabular}{ccccc}
\toprule
\textbf{Seed} & \textbf{Test loss} & \textbf{IIA}$_{k{=}2}$ &
\textbf{Equiv.} & \textbf{Mean $d_\Gr$ to others} \\
\midrule
42   & 0.000 & 1.000 & 0.995 & 2.08 \\
137  & 0.000 & 0.980 & 0.991 & 2.11 \\
256  & 0.000 & 1.000 & 0.998 & 2.05 \\
512  & 0.000 & 0.960 & 0.989 & 2.09 \\
1024 & 0.000 & 1.000 & 0.997 & 2.07 \\
2024 & 0.000 & 0.940 & 0.985 & 2.13 \\
3000 & 0.000 & 1.000 & 0.993 & 2.06 \\
4096 & 0.000 & 0.980 & 0.990 & 2.10 \\
5555 & 0.000 & 1.000 & 0.996 & 2.08 \\
9999 & 0.000 & 1.000 & 0.992 & 2.09 \\
\midrule
\multicolumn{2}{l}{Mean $\pm$ std} & $0.986 \pm 0.020$ & $0.993 \pm 0.004$ &
$2.09 \pm 0.02$ \\
\bottomrule
\end{tabular}
\caption{Per-seed results for modular multiplication ($p = 113$, 1-layer).
All 10 seeds grok and achieve high IIA and equivariance, but their DAS
subspaces are near-orthogonal (mean pairwise $d_\Gr = 2.09$, compared to
the random-subspace expectation of $\approx 2.20$ on $\Gr(2, 128)$).}
\label{tab:seed_data}
\end{table}


\section{Pi-SAE Ablation Details}
\label{app:ablation}

The $2 \times 2$ ablation (Table~\ref{tab:ablation}) tests whether both
components of the Structured pi-SAE---the label-conditional prior and the
sparsity constraint---are individually necessary. This appendix provides
the full per-operation breakdown including reconstruction MSE, which
serves as a built-in falsifiability diagnostic.

\vspace{1em}
\begin{table}[H]
\centering
\small
\setlength{\tabcolsep}{3.5pt}
\renewcommand{\arraystretch}{1.15}
\begin{tabular}{@{}l|cc|cc|cc|cc@{}}
\toprule
& \multicolumn{2}{c|}{\textbf{VAE}} & \multicolumn{2}{c|}{\textbf{SAE}}
& \multicolumn{2}{c|}{\textbf{pi-VAE}} & \multicolumn{2}{c}{\textbf{Str.\ pi-SAE}} \\
\textbf{Operation} & IIA & MSE & IIA & MSE & IIA & MSE & IIA & MSE \\
\midrule
\multicolumn{9}{l}{\textit{Grokked operations}} \\
Addition        & 0.02 & 0.8 & 1.00 & 1.1  & 0.00 & 0.7  & \textbf{1.00} & 0.6 \\
Multiplication  & 0.00 & 0.9 & 0.72 & 1.3  & 0.00 & 0.8  & \textbf{1.00} & 0.7 \\
Quartic sum     & 0.02 & 1.0 & 0.32 & 1.5  & 0.00 & 0.9  & \textbf{1.00} & 0.8 \\
IOI (GPT-2)     & 0.56 & 2.1 & 0.83 & 3.4  & 0.95 & 1.8  & \textbf{0.98} & 1.2 \\
\midrule
\multicolumn{9}{l}{\textit{Non-grokked operations}} \\
Mixed product   & 0.03 & 18  & 0.01 & 21   & 0.01 & 16   & 0.04 & 11 \\
Squaring        & 0.00 & 22  & 0.00 & 24   & 0.00 & 19   & 0.00 & 14 \\
\bottomrule
\end{tabular}
\renewcommand{\arraystretch}{1.0}
\caption{Full $2 \times 2$ ablation with reconstruction MSE. Grokked operations
have MSE $\leq 1.2$; non-grokked operations have MSE $\geq 11$---an order of
magnitude gap that serves as a falsifiability diagnostic. High MSE with
near-zero IIA means the model has no structured causal variable to recover.
The plain SAE degrades on harder operations (multiplication 0.72, quartic sum
0.32) because without the structured prior, the sparse encoder has no target
direction. The pi-VAE achieves IIA $\approx 0$ despite correct prior direction
because without sparsity, the encoder distributes the causal signal across all
dimensions.}
\label{tab:ablation_full}
\end{table}

\paragraph{Control experiments.}
Five controls validate that the Structured pi-SAE's success requires
correct supervision and the causal-nuisance partition, ruling out
trivial explanations:

\vspace{0.5em}
\begin{table}[H]
\centering
\small
\begin{tabular}{lcc}
\toprule
\textbf{Control} & \textbf{IIA} & \textbf{Purpose} \\
\midrule
Random-label pi-SAE    & $\leq 0.02$  & Labels carry no causal signal \\
Reconstruction-only ($\alpha{=}0$) & $\leq 0.005$ & No supervision to pin direction \\
Untrained pi-SAE       & 0.000        & Architecture alone is insufficient \\
Plain VAE (no split)   & $\leq 0.08$  & Causal-nuisance partition is necessary \\
Random subspace (DAS)  & $\leq 0.01$  & Baseline for chance-level IIA \\
\bottomrule
\end{tabular}
\caption{Control experiments on grokked multiplication ($p = 113$).
All controls achieve near-zero IIA, confirming that the Structured pi-SAE's
success requires correct labels, supervision, and the causal-nuisance split.}
\label{tab:controls}
\end{table}
